\section{Revisão de literatura e trabalhos relacionados}
\label{sec:revisao}

\subsection{Protocolo de busca e seleção}

Será conduzido um mapeamento da literatura conforme as recomendações atualizadas
de \cite{petersenEtAl2015}. O protocolo será congelado antes da seleção final e
registrará bases, strings, filtros, datas, duplicatas, decisões por título e
resumo, leitura completa, avaliação de qualidade e extração. As bases definidas
são IEEE Xplore, ACM Digital Library, Scopus e SciELO. Google Scholar será usado
apenas de forma complementar e para encadeamento de referências.

As strings combinam três grupos: (i) \textit{misfueling}, \textit{wrong fuel} e
\textit{fueling error}; (ii) prevenção, intertravamento, autorização, ergonomia
e fatores humanos; e (iii) identificação veicular, placa, RFID e dispensador.
Serão incluídos estudos revisados por pares, em português, inglês ou espanhol,
com texto completo e relação direta com ao menos uma questão da revisão.
Patentes, publicidade, notícias e manuais serão excluídos da síntese científica,
mas poderão ser analisados separadamente como documentação técnica.

Cada estudo receberá avaliação qualitativa quanto à clareza do objetivo,
descrição de contexto e amostra, reprodutibilidade do método, coerência das
métricas, separação entre resultados e interpretação e discussão de limitações.
O protocolo completo está versionado em
\texttt{research/protocol/revisao-literatura.md}; o diário e a matriz de extração
ficam em \texttt{research/literature/}.

\pendente{executar as buscas em todas as bases e inserir o fluxograma de seleção
com contagens verificáveis.}

\subsection{Síntese preliminar}

A Tabela~\ref{tab:trabalhos-relacionados} reúne as fontes centrais já
verificadas. Ela não representa o resultado final do mapeamento.

\begin{table}[H]
  \centering
  \caption{Trabalhos e documentos relacionados já verificados}
  \label{tab:trabalhos-relacionados}
  \small
  \begin{tabularx}{\textwidth}{@{}L{3.2cm}L{3.3cm}Y@{}}
    \toprule
    \textbf{Fonte} & \textbf{Método ou artefato} & \textbf{Relação com esta pesquisa} \\
    \midrule
    Adams e David (2007) & Análise hierárquica da tarefa e dados de 23 pessoas que cometeram o erro & Evidencia a combinação de fatores humanos e ambientais e recomenda modificar leiaute ou equipamento. \\
    Mikkonen e Engman (2013) & Análise técnica das consequências da troca entre gasolina e diesel & Sustenta que pequenas quantidades de gasolina podem degradar propriedades relevantes do diesel e que a gravidade depende do cenário. \\
    Taylor (1978) & Levantamento em postos de grandes cidades canadenses & Oferece registro histórico de incidência e de alterações em bocais, mas não permite inferir frequência atual ou brasileira. \\
    Gilbarco (manual técnico) & Identificação automática por RFID e autorização da bomba & Documenta uma abordagem comercial distinta. Não é evidência científica de eficácia nem define a arquitetura deste projeto, que não comanda bombeamento. \\
    \bottomrule
  \end{tabularx}
  \\[2pt]\fontepropria
\end{table}

O estudo de \cite{adamsDavid2007} é o mais diretamente relacionado à origem do
problema: trata o abastecimento como uma sequência suscetível a condições que
favorecem o erro. Já \cite{mikkonenEngman2013} concentra-se nas consequências
técnicas da contaminação entre gasolina e diesel, o que ajuda a justificar a
atuação preventiva, mas não avalia uma barreira automatizada. A documentação
\cite{gilbarcoFuelPoint2018} mostra que controladores comerciais já realizam
autorização eletrônica de abastecimento por identificação automática. A solução
proposta neste TCC difere ao investigar consulta veicular por placa no contexto
regulatório brasileiro, política explícita de compatibilidade, comportamento
\textit{fail-closed} e análise econômica para uma rede específica.

Até esta etapa, não foi identificado estudo revisado por pares que avalie
conjuntamente esses quatro elementos. Essa observação é um \textbf{achado
preliminar}, limitado às fontes-semente, e somente poderá ser apresentada como
lacuna consolidada após a execução integral do protocolo.
